\documentclass[10pt,a4paper]{amsart}
\usepackage[margin=19mm]{geometry}
\usepackage{amsmath,amssymb,mathtools}
\usepackage[scr=rsfs,cal=euler]{mathalfa}
\usepackage{xfrac}
\usepackage[unicode,hidelinks,bookmarks=false]{hyperref}
\newcommand{\ie}{i.e.\ }
\newcommand{\cf}{cf.\ }
\newcommand{\resp}{respectively\ }
\newcommand{\Subsection}{\subsection}
\providecommand{\nobreakdash}{\nobreak\hbox{-}}
\setlength{\emergencystretch}{2em}
\multlinegap0pt
\newcommand{\bCP}{\mathbb{CP}^n}
\usepackage{tikz}
\usetikzlibrary{cd}
\usepackage{amssymb}
\usepackage{enumerate}
\usepackage[shortlabels]{enumitem}
\newtheorem{defn}{Definition}
\newtheorem{prop}{Proposition}
\newcommand{\bC}{\mathbb{C}}
\newcommand{\bR}{\mathbb{R}}
\title{Zoll Manifolds of Type $\bCP$ with Entire Grauert Tubes}
\author{Chi Li, Kyobeom Song}
\date{Source: arXiv:2501.15682v2 (CC BY 4.0)}
\begin{document}
\maketitle

\noindent\emph{Source and licence.} Zoll Manifolds of Type $\bCP$ with Entire Grauert Tubes. Version arXiv:2501.15682v2 (CC BY 4.0), distributed by arXiv under the Creative Commons Attribution 4.0 International licence, http://creativecommons.org/licenses/by/4.0/, as recorded in the arXiv licence metadata for this exact version and shown on the abstract page https://arxiv.org/abs/2501.15682v2. Full licence notice: \texttt{licenses/arxiv-2501.15682-CC-BY-4.0.txt}.\par\smallskip

\noindent\textit{Editorial scope.} Counted unit: the article's prop bundle (source0.tex lines 401--407). The article's complete original proof of it is delivered with it (source0.tex lines 409--452, one \texttt{proof} environment of 44 lines). No mathematics is written, repaired or re-derived: the statement and the proof are the article's own lines, in the article's own notation, and no retained line is substituted or re-flowed.  Retained uncounted as the source-local context the proof needs: lines 370--397.  Nothing was omitted from any retained span, so the package carries no omission record.  The retained lines cite 1 works, whose entries are transcribed verbatim from the article's own bibliography source in the e-print and delivered with short editorial labels recorded in the item's reference mapping.  The retained lines contain the article's own diagram code, which the wrapper typesets with the standard tikz packages; no image file is delivered and the package declares no asset dependency.  Editorial formatting only, in addition: an AMS \texttt{amsart} preamble that restates the article's own declarations and macros used by the retained lines, namely the environments \texttt{defn}, \texttt{prop} on separate counters, together with the standard packages those lines need.  The retained environments print in this document as Definition 1, Proposition 1 in order of appearance, whereas the article prints the same items with the numbering of its own full document.  The complete native source of the article as distributed in the arXiv e-print for this exact version is bundled unchanged as \texttt{sources/172206/source0.tex}.
\begin{defn}\label{chart}
Let $D$ be the set of oriented images of geodesics in $M$. Let $\gamma$ be a unit-speed geodesic in $M$ with $\gamma(0)=p$ and $\gamma'(0)=v$. An element of $D$ corresponding to $\gamma$ is denoted $0_\gamma$.
Consider the normal vector space $N_vM=\gamma'(0)^\perp \subset T_pM$.
For $u,w\in N_v M$, define a geodesic $\gamma(u,w)$ by
\[
\gamma(u,w)(s):=\exp_{\exp_p u}\!\Bigl( s\, P_u\!\Bigl(\frac{v+w}{\sqrt{1+\|w\|^2}}\Bigr)\Bigr),
\]
where $P_u$ denotes parallel transport along the curve $t\mapsto \exp_p(tu)$ from $p$ to $\exp_p u$. For a sufficiently small open neighborhood $U\subset N_v M$ of $0\in N_v M$, we define a smooth local parametrization $\varphi'_\gamma$ of $D$ around $0_\gamma$ by
\[
\varphi'_\gamma \colon U\times U \subset N_vM\times N_vM \cong \bR^{2m-2} \longrightarrow D,
\qquad
\varphi'_\gamma(u,w):=0_{\gamma(u,w)}.
\]
Let $\Delta:=\{z\in \bC \mid |z|<1\}$. Define $\varphi_\gamma \colon U\times U\times \Delta \longrightarrow X$ by
\[
\varphi_\gamma\!\bigl(u,w,re^{i\theta}\bigr)=
\begin{cases}
\dfrac{l}{4\pi}\cosh^{-1}\left(\dfrac{1}{r^2}\right)\,\gamma(u,w)'\!\left(\dfrac{l}{2\pi}\theta\right), & r\in (0,1),\\[0.4em]
0_{\gamma(u,w)}\in D, & r=0.
\end{cases}
\]
Finally, define $\pi \colon X\setminus M \to D$ by
\[
\pi\bigl(a\,\gamma'(b)\bigr)=0_\gamma \quad \text{for any } a,b>0,
\qquad\text{and}\qquad
\pi(0_\gamma)=0_\gamma.
\]
\end{defn}

\begin{prop}\label{bundle}
Let $M$ be a Zoll manifold, and consider a set $X := TM \cup D$ and $\varphi_\gamma$, $\pi$ defined in Definition \ref{chart}.
\begin{enumerate}
\item The set $X$ admits a unique smooth structure extending the given smooth structure on $TM$ such that each map $\varphi_\gamma$ is a smooth local parametrization.
\item $(X\setminus M, D, \pi, \Delta)$ defines the structure of an open disc bundle. In particular, $X\setminus M$ deformation retracts onto $D$.
\end{enumerate}
\end{prop}

\begin{proof}
Let $\alpha,\beta$ be unit-speed geodesics in $M$ with $\alpha(0)=p$, $\alpha'(0)=v$, and $\beta(0)=p'$, $\beta'(0)=v'$. Consider the maps $\varphi'_\alpha,~\varphi_\alpha,~\varphi'_\beta,~\varphi_\beta$ constructed as in Definition~\ref{chart}. By definition, on the overlap we have
\[
\varphi_\beta^{-1}\circ \varphi_\alpha\!\bigl(u,w,re^{i\theta}\bigr)
=
\Bigl((\varphi_\beta')^{-1}\circ \varphi_\alpha'(u,w),\ re^{i(\theta+t)}\Bigr),
\]
where $t=t(u,w)$ is a function such that for $(\varphi_\beta')^{-1}\circ \varphi_\alpha'(u,w)=(u',w')$,
\[
\alpha(u,w)(0)=\beta(u',w')\left(\dfrac{l}{2\pi}t\right). 
\]
Since each ingredient in this defining relation depends smoothly on $(u,w)$, it follows that $t(u,w)$ is a smooth function.

Writing $z=re^{i\theta}$, we have $re^{i(\theta+t)}=z\,e^{it}$, hence
\[
\varphi_\beta^{-1}\circ \varphi_\alpha(u,w,z)
=
\Bigl((\varphi_\beta')^{-1}\circ \varphi_\alpha'(u,w),\ z\,e^{i t(u,w)}\Bigr),
\]
which is smooth (in particular, smooth at $z=0$). Moreover, the map $\varphi_\gamma$ restricts to a smooth map
\[
\varphi_\gamma\big|_{U\times U\times (\Delta\setminus\{0\})}\colon U\times U\times (\Delta\setminus\{0\})\longrightarrow TM
\]
with respect to the standard smooth structure on $TM$. Hence, for each point of $M\subset TM$, we may take the usual local charts of $TM$ around that point; together with the local parametrizations $\varphi_\gamma$, these charts cover $X$, and all transition maps between them are smooth. Therefore, they determine a unique smooth manifold structure on $X$ by the smooth manifold chart lemma (cf.~\cite[Lemma~1.35]{L03}).

Finally, since $\pi\circ\varphi_\gamma(u,w,z)=0_{\gamma(u,w)}=\varphi'_\gamma(u,w)$, the following diagram commutes:
\[
\begin{tikzcd}[column sep=huge]
\varphi_\gamma(U\times U\times \Delta)
  \arrow[r, "(\varphi'_\gamma\times \mathrm{id}_{\Delta})\circ \varphi_\gamma^{-1}"]
  \arrow[d, "\pi"']
&
\varphi'_\gamma(U\times U)\times \Delta
  \arrow[dl, "\mathrm{proj}_{\varphi'_\gamma(U\times U)}"]
\\
\varphi'_\gamma(U\times U) &
\end{tikzcd}
\]
Therefore, the pairs
\[
\Bigl(\varphi_\gamma(U\times U\times \Delta),\ (\varphi'_\gamma\times \mathrm{id}_{\Delta})\circ \varphi_\gamma^{-1}\Bigr)
\]
form a local trivialization of the disc bundle $(X\setminus M, D, \pi, \Delta)$, proving the second statement. The deformation retraction onto $D$ follows by radial contraction in each disc fiber.
\end{proof}

\begin{thebibliography}{99}
\bibitem[L03]{L03} Lee, J.~M. \emph{Introduction to Smooth Manifolds}, Graduate Texts in Mathematics, Vol.~218, Springer, New York, 2003.
\end{thebibliography}
\end{document}
